\section{A comparison tree and cost-weighted curvature}\label{sec:comparison}

Multiscale restrictions of alignment paths, controlled by bounds on their
total deviation, are central to Mao and
Rubinstein's approximation scheme~\cite[Sections~2.3 and~6.1]{MaoRubinstein2026}.
Here the weights are warmup cost estimates. They need not add exactly from
children to parent, so our energy calculation also charges their mass
discrepancy and the difference between a comparison action and a minimum
wide action.

Fix a seed table satisfying the hypotheses of Theorem~\ref{thm:refinement}.
We construct a comparison tree from exact alignments.  Its states are fixed
before the randomness of this refinement is drawn.  The warmup values supply
weights on this tree.  A telescoping energy calculation controls the total
cost of nodes with large curvature or a large gap between the comparison
action and the minimum wide action.  At the other nodes, the velocity test
has two useful outcomes.  Passing it places the comparison action in the
computed band; failing it means that length imbalance nearly accounts for
the warmup cost.  Section~\ref{sec:regret} will use these alternatives to
bound the refinement error.

\subsection{Exact paths and rounded child states}

Start at a state $(I,q)$ of true cost $D=c(I,q)$.  If $D=0$, the good-seed
bounds force all its estimates to be zero, so assume $D>0$.  Stop a branch at
a leaf of the interval tree or at a state of cost zero.  At each remaining
node, write $c=c(I,q)>0$, choose the dyadic scale $b$ with $b\le c<2b$, and put
\[
 r=r_b(q),\qquad e=|q-r|_1\le 2b/P.
\]
Choose an optimal alignment for $(I,r)$ and points on this path at the child
cuts, in their path order.  Use its exact endpoints at the outside cuts.
For repeated cuts, repeated path points are permitted; alternatively a
vertical segment can be assigned to one of the intervening empty children.
Round the internal target anchors down to the grid at scale $b$.  The resulting
child states $q_1,\ldots,q_M$ remain ordered and share endpoints, with outside
endpoints exactly $r$.  Denote this action by $\alpha^*$, and continue the
construction at its child states.

Fix a deterministic order for all choices of optimal paths and cut points.
The entire comparison tree, including its stopping states, then depends only
on the input strings and its starting state.  In particular it is independent
of every group sample and online draw in the refinement.  This tree is used
only in the proof.

\begin{lemma}[Comparison costs]\label{lem:comparison-tree}
At each nonterminal of the comparison tree, $(b,r)$ is a representative,
$b$ satisfies the local scale restriction~\eqref{eq:local-scales}, and
$\alpha^*$ is a wide action whose endpoint shifts are within $6b$ of
$p_0=r^- -l$.  Writing $c_i=c(I_i,q_i)$, we have
\begin{equation}\label{eq:comparison-split}
 c(I,r)\le\sum_i c_i\le (1+4M/P)c,
 \qquad c\le Le+\sum_i c_i.
\end{equation}
Let $C_{\mathrm{term}}$ be the sum of terminal true costs.  Then
\begin{equation}\label{eq:comparison-totals}
 \begin{aligned}
 C_{\mathrm{term}}&\le(1+8MJ/P)D\le2D,\\
 \sum_{\mathrm{nonterminal}}c&\le2JD,
 &\sum_{\mathrm{nonterminal}}e&\le4JD/P.
 \end{aligned}
\end{equation}
The sum of true costs at any depth is at most $2D$, also if stopped costs are
carried unchanged to subsequent depths.  At every terminal, $m=c$ and every
refinement estimate equals $c$.
\end{lemma}

\begin{proof}
If the grid spacing is one, no rounding changes any endpoint.  Otherwise an
endpoint moves by less than $b/P$.  Thus
\[
 c-e\le c(I,r)\le c+e.
\]
The initial-cone bounds give
\[
 w_0(I,r)\ge c(I,r)/a\ge b/(4a),
 \qquad
 w_0(I,r)\le2F c(I,r)\le8Fb,
\]
so $(b,r)$ is a representative.  At $q$, the same bounds imply
$w_0(I,q)/(16F)\le b\le8a w_0(I,q)$, as required for
\eqref{eq:local-scales}.

Along an alignment, each change of shift is caused by an insertion or
deletion and is at most the cost of that part of the path.  Hence every
unrounded anchor shift differs from $p_0$ by at most $c(I,r)$.  Rounding adds
less than $b/P$, so the distance is at most
\[
 c+e+b/P<(2+3/P)b<6b.
\]
This proves the wide-action and radius assertions.

The exact subpath costs before rounding sum to $c(I,r)$.  Each internal anchor
affects the endpoints of two child problems.  Endpoint Lipschitz continuity
therefore gives
\[
 \sum_i c_i\le c(I,r)+2(M-1)b/P
 \le c+2Mb/P\le(1+4M/P)c.
\]
Concatenation gives the other inequality in the first part of
\eqref{eq:comparison-split}; it also gives
$c\le c(I,r)+e\le\sum_i c_i+Le$.
These arguments apply without change to repeated cuts and empty child
intervals.

At each successive frontier of the tree, carry terminal costs unchanged.
The frontier sum grows by a factor at most $1+4M/P$ in one step.  The parameter
bounds give $4MJ/P\le1/2$, and consequently
\[
 (1+4M/P)^J\le1+8MJ/P\le2.
\]
This proves the terminal and level bounds.  There are at most $J$
nonterminal levels.  Summing $e\le2b/P\le2c/P$ proves the remaining bounds
in~\eqref{eq:comparison-totals}.

At a leaf, every table uses the exact answer.  At a zero-cost state, the
good seed is zero and the upper bound $w_0\le2Fc$ forces all warmup and
refinement values to be zero.  These are precisely the terminal cases.
\end{proof}

\subsection{Mass and energy telescopes}

If the warmup masses added exactly from children to parent, the energies
$z^2/m$ would telescope up to endpoint-rounding errors.  The masses do not
add exactly: a parent's mass can be either
larger or smaller than its children's total.  The warmup inequalities bound
the total upward discrepancy, and conservation of mass over the comparison
tree then bounds the total downward discrepancy as well.  This is why the
energy argument can use cost estimates instead of exact additive weights.

At a nonterminal use the representative and child states just constructed,
and define
\[
 m_q=m(I,q),\qquad m_i=m(I_i,q_i),\qquad W=\sum_i m_i,
 \qquad G=W-Z,\qquad R=m_q-W.
\]
Here $Z$ is the minimum wide-action cost with child table $m$, as in
Section~\ref{sec:bands}.  The comparison action has zero connection cost,
so $W\ge Z$.  The lower invariants and the comparison costs imply
\begin{equation}\label{eq:comparison-masses}
 0<\max\{c(I,r)/a,|z(I,r)|\}\le Z\le W\le4Fc<8Fb.
\end{equation}
In particular, the denominator $Z$ defining $\widehat v$ is positive at
every nonterminal comparison representative.

Write $z_i=z(I_i,q_i)$ and $z_r=z(I,r)$.  Since the child targets share
endpoints, $\sum_i z_i=z_r$.  Define the curvature quantity
\begin{equation}\label{eq:curvature-definition}
 K=\sum_i\frac{z_i^2}{m_i}-\frac{z_r^2}{W},
\end{equation}
with $0/0=0$.  A zero $m_i$ necessarily has $z_i=0$ by the lower invariant.
For $v=z_r/W$, the weighted variance identity is
\begin{equation}\label{eq:weighted-variance}
 K=\sum_{i:m_i>0}\frac{(z_i-vm_i)^2}{m_i}\ge0.
\end{equation}
These weights are costs supplied by the warmup; they are not interval
lengths.

\begin{lemma}[Total curvature]\label{lem:curvature-telescope}
On the fixed comparison tree, put $h=Le+w_0(I,q)/T$ at each nonterminal.
Then
\begin{equation}\label{eq:curvature-structural}
 \sum_{\mathrm{nonterminal}}(K+G)
 \le4C_{\mathrm{term}}+3\sum_{\mathrm{nonterminal}}e
                       +4\sum_{\mathrm{nonterminal}}h.
\end{equation}
Consequently,
\begin{equation}\label{eq:curvature-total}
 \sum_{\mathrm{nonterminal}}(K+G)
 \le5000F(1+J/P+J/T)D.
\end{equation}
\end{lemma}

\begin{proof}
First compare the parent and child warmup averages.  For each wide action,
average the parent warmup inequalities over indices $1,\ldots,T$, using
child indices $0,\ldots,T-1$.  Minimize the resulting action cost.  This gives
\[
 m_q-\frac{w_0(I,q)-w_T(I,q)}T\le Le+Z.
\]
Thus
\begin{equation}\label{eq:warmup-mass-slack}
 R+G\le h,\qquad R_+\le h,\qquad G\le R_-+h,
 \qquad R_\pm=\max\{\pm R,0\}.
\end{equation}
The masses cancel exactly across the tree:
\[
 \sum R=m_{\mathrm{root}}-C_{\mathrm{term}},
 \qquad
 \sum R_-\le\sum R_++C_{\mathrm{term}}
             \le\sum h+C_{\mathrm{term}}.
\]
In particular, summing $G\le R_-+h$ gives
$\sum G\le C_{\mathrm{term}}+2\sum h$.

Next define the energy of any comparison state by
\[
 E(I,q)=z(I,q)^2/m(I,q),
\]
again using $0/0=0$.  At a nonterminal put $z_q=z(I,q)$.  We have
$|z_q|\le m_q$, $|z_r-z_q|\le e$, and $|z_r|\le W$.  Moreover,
\[
 e\le2c/P\le2a m_q/P\le m_q,
\]
because $a\le A_{\rm init}$ and $2A_{\rm init}<P$.  Therefore
$|z_r|\le2m_q$.  The identity
\[
 \frac{z_q^2}{m_q}-\frac{z_r^2}{W}
 =\frac{z_q^2-z_r^2}{m_q}
  +\frac{z_r^2(W-m_q)}{m_qW}
\]
has first term at most $3e$.  If $W\le m_q$, its second term is nonpositive;
otherwise the bounds $|z_r|\le W$ and $|z_r|\le2m_q$ bound it by
$2(W-m_q)=2R_-$.  It follows that
\begin{equation}\label{eq:local-energy-telescope}
 K\le\sum_iE(I_i,q_i)-E(I,q)+3e+2R_-.
\end{equation}
No zero denominator occurs at a nonterminal, by
\eqref{eq:comparison-masses}; zero child masses have already been treated.

Terminal energies are at most their true costs, and the root energy is
nonnegative.  Summing~\eqref{eq:local-energy-telescope} therefore gives
\[
 \sum K\le C_{\mathrm{term}}+3\sum e+2\sum R_-
          \le3C_{\mathrm{term}}+3\sum e+2\sum h.
\]
Adding the preceding bound on $\sum G$ proves
\eqref{eq:curvature-structural}.  It remains only to substitute the
comparison-tree bounds.  By Lemma~\ref{lem:comparison-tree} and $w_0\le2Fc$,
\begin{equation}\label{eq:total-warmup-slack}
 \sum h\le512FJD/P+4FJD/T.
\end{equation}
Together with~\eqref{eq:comparison-totals}, this yields
\[
 \sum(K+G)\le8D+(2048F+12)JD/P+16FJD/T,
\]
which implies~\eqref{eq:curvature-total}.
\end{proof}

Call a nonterminal \emph{bad} if $K>bH^{-1/4}$ or $G>bH^{-1/4}$; otherwise
call it \emph{non-bad}.  These are conditions on the deterministic comparison
tree and its warmup values.

\begin{corollary}\label{cor:bad-node-cost}
The total true cost of bad nodes satisfies
\begin{equation}\label{eq:bad-node-cost}
 \sum_{\mathrm{bad}}c
 \le10000FH^{1/4}(1+J/P+J/T)D.
\end{equation}
\end{corollary}

\begin{proof}
At a bad node, $c<2b<2H^{1/4}(K+G)$.  Sum and apply
Lemma~\ref{lem:curvature-telescope}.
\end{proof}

\subsection{Coverage by the computed bands}

The preceding estimate controls the cost of exceptional nodes.  We now show
that at every remaining node which passes the velocity test, the band and
rounding guards retain the comparison action.  This is the point where
cost-weighted curvature produces a usable prediction for the algorithm.

\begin{lemma}[Comparison-band coverage]\label{lem:comparison-band}
At a non-bad comparison node with $|\widehat v|\le1-\tau$, the action
$\alpha^*$ belongs to $\mathcal C_{b,r}$.  Its representative is not skipped:
the band is nonempty, $h_\Sigma\le64Fb$, and the rounded action $\beta(b,r)$
passes both the wide-action and imbalance guards.
\end{lemma}

\begin{proof}
Let $p_j^*$ be the comparison shift at internal cut $j$, and write
$M_j=\sum_{i\le j}m_i$.  The comparison prefix is feasible for
$\Phi_j(p_j^*)$, so this prefix cost is at most $M_j$.  Conversely, concatenate
any feasible prefix in that minimum with the comparison suffix.  The prefix's
final jump is exactly its connection to the suffix.  The resulting wide
action has cost at least $Z$, and hence
\[
 M_j-G\le\Phi_j(p_j^*)\le M_j.
\]
With $v=z_r/W$, we also have
\[
 |v-\widehat v|W=|z_r|G/Z\le G.
\]
Weighted Cauchy--Schwarz in~\eqref{eq:weighted-variance} gives
\[
 \sum_i|z_i-vm_i|\le\sqrt{KW}.
\]
The shared endpoints imply $p_j^*-p_0=\sum_{i\le j}z_i$.  Combining the last
three observations yields
\begin{equation}\label{eq:band-coverage}
 |p_j^*-p_0-\widehat v\Phi_j(p_j^*)|
 \le\sqrt{KW}+2G\le\tau\theta b/M.
\end{equation}
For completeness, the last numerical inequality follows from
$K,G\le bH^{-1/4}$, $W<8Fb$, and the parameter bounds:
\[
 \frac M b(\sqrt{KW}+2G)
 \le\sqrt{8\ell}\,H^{-3/40}+2H^{-1/5}
 \le\ell^{-44}=\tau\theta.
\]
Indeed $\ell^{200}\le H^{1/100}$ bounds the middle expression by
$\sqrt{8\ell}\,\ell^{-1500}+2\ell^{-4000}$.
The radius and shared-endpoint conditions were proved in
Lemma~\ref{lem:comparison-tree}.  Thus~\eqref{eq:band-coverage} proves
$\alpha^*\in\mathcal C_{b,r}$.

The diameter bound for the band implies that any of its child states is
within $4\theta b/M$ of $q_i$ in endpoint metric.  Using the
$L_0$-Lipschitz property of $w_0$ and the factor of at most two from dyadic
ceiling, we get
\[
 \begin{aligned}
 h_\Sigma
 &\le2b+2\sum_i w_0(I_i,q_i)+8L_0\theta b\\
 &\le2b+4F\sum_i c_i+8L_0\theta b
 \le64Fb.
 \end{aligned}
\]
The last step uses~\eqref{eq:comparison-split}, $c<2b$, and the parameter
range.  This verifies the sum guard.

Under the fixed good seed, the rounding verification following
\eqref{eq:rounded-guard} applies as soon as this sum guard holds.
It shows that $\beta(b,r)$ is a wide action and satisfies the imbalance
guard.  Rounding preserves the order of each child's two endpoints;
different child spacings may separate shared endpoints, with total
connection charge at most $2\theta h_\Sigma$.  Thus every guard passes.
\end{proof}
